\section{SPECTRA-DA}
\label{sec:method}

\subsection{Problem setting and sealed selection}

We study model selection after UGDA training. Let
$\mathcal G_S=(A_S,X_S,Y_S)$ be a labeled source graph and
$\mathcal G_T=(A_T,X_T)$ an unlabeled target graph. A frozen candidate bank
$\mathcal H=\{h_m\}_{m=1}^M$ contains different adaptation algorithms,
hyperparameters, random seeds, and checkpoints. The selector may access source
labels, both public graphs, and candidate probabilities $P_m$ and embeddings
$Z_m$, but never target labels $Y_T$. It outputs
\begin{equation}
  \widehat m=\arg\min_{m\in[M]}\widehat R_T(h_m).
  \label{eq:selection}
\end{equation}
For single-label node classification, Micro-F1 equals accuracy, so true risk is
$R_T(h_m)=1-\operatorname{MicroF1}_T(h_m)$. Evaluation measures the regret of
$h_{\widehat m}$ relative to the best candidate in the same immutable bank.

\benchmark enforces this information boundary operationally. Public trajectory
artifacts contain source-validation and unlabeled target outputs but no target
label field. Candidate-bank hashes bind every score vector to a fixed ordering.
The selector process cannot import the sealed evaluator, and access counters
record attempted label reads. Only after code, candidates, and scores are
frozen does an isolated evaluator load $Y_T$ to compute aggregate metrics. This
protocol prevents direct target-oracle selection and repeated feedback from a
hidden target leaderboard.

\subsection{Tight spectral risk decomposition}

Global disagreement cannot distinguish graph-smooth errors from locally
oscillatory errors. We therefore decompose predictions over the target graph,
following the graph-signal-processing view of Laplacian spectral filters
~\citep{shuman2013emerging,hammond2011wavelets}.
Let
\begin{equation}
  L_T=I-D_T^{-1/2}A_TD_T^{-1/2}
\end{equation}
be the symmetric normalized Laplacian. We define $B$ Gaussian spectral windows
$\widetilde g_b$ and normalize them pointwise,
\begin{equation}
  \widetilde g_b(\lambda)
  =\exp\!\left(-\frac{(\lambda-\mu_b)^2}{2\sigma^2}\right),
  \qquad
  g_b(\lambda)
  =\frac{\widetilde g_b(\lambda)}
  {\sqrt{\sum_j\widetilde g_j(\lambda)^2}}.
  \label{eq:tight_windows}
\end{equation}
The resulting filters $G_b=g_b(L_T)$ satisfy
$\sum_bG_b^\top G_b=I$. A Chebyshev approximation applies all filters jointly
without eigendecomposition~\citep{defferrard2016convolutional}; the implementation rejects configurations whose
maximum frame error exceeds the frozen tolerance.

Let $H_m\in\{0,1\}^{n_T\times C}$ be the hard one-hot prediction of candidate
$m$, and define the unknown error signal $E_m=H_m-Y_T$. Its band risk is
\begin{equation}
  r_{m,b}=\frac{1}{n_T}\lVert G_bE_m\rVert_F^2.
  \label{eq:band_risk}
\end{equation}

\begin{lemma}[Exact hard-error decomposition]
If the spectral filter bank is tight, then
\begin{equation}
  \operatorname{Err}_T(h_m)=\frac{1}{2}\sum_{b=1}^{B}r_{m,b}.
  \label{eq:exact_decomposition}
\end{equation}
\end{lemma}
\begin{proof}
Tightness gives
$\sum_b r_{m,b}=n_T^{-1}\operatorname{tr}
(E_m^\top(\sum_bG_b^\top G_b)E_m)
=n_T^{-1}\lVert E_m\rVert_F^2$.
Two one-hot class vectors have squared distance zero when correct and two when
incorrect, hence $\lVert E_m\rVert_F^2=2n_T\operatorname{Err}_T(h_m)$.
\end{proof}

The same frame identity decomposes Brier risk for soft probabilities $P_m$,
but not classification error exactly. We use hard predictions for the theorem
and soft predictions as a smoother empirical signal. An approximate-frame
bound is given in Appendix~\ref{app:frame_error}.

\subsection{Recovering individual risk from pair disagreement}

Although $r_{m,b}$ is hidden, pairwise spectral disagreement is observable:
\begin{equation}
  d_{mn,b}=\frac{1}{n_T}
  \lVert G_b(H_m-H_n)\rVert_F^2.
  \label{eq:disagreement}
\end{equation}
Since $H_m-H_n=E_m-E_n$, expansion yields the exact relation
\begin{equation}
  d_{mn,b}=r_{m,b}+r_{n,b}-2c_{mn,b},
  \qquad
  c_{mn,b}=\frac{1}{n_T}
  \langle G_bE_m,G_bE_n\rangle_F.
  \label{eq:pair_identity}
\end{equation}
Thus disagreement fails to identify individual risk only through the
cross-model error covariance $c_{mn,b}$.

Let $A_{\mathrm{pair}}\in\mathbb R^{\binom M2\times M}$ have one row
$e_m+e_n$ per unordered candidate pair. Stacking Equation~\eqref{eq:pair_identity}
within a band gives
\begin{equation}
  \mathbf d_b=A_{\mathrm{pair}}\mathbf r_b-2\mathbf c_b.
  \label{eq:pair_system}
\end{equation}
For zero covariance, the untruncated complete-pair solution has the closed
form
\begin{equation}
  \widetilde r_{m,b}
  =\frac{s_{m,b}-S_b/(M-1)}{M-2},
\end{equation}
where $s_{m,b}=\sum_{n\neq m}d_{mn,b}$ and
$S_b=\sum_{i<j}d_{ij,b}$. The practical estimator uses nonnegative least
squares because finite-sample and covariance errors can otherwise yield
negative risks.

\begin{theorem}[Covariance-controlled recovery]
For a complete pair graph,
$A_{\mathrm{pair}}^\top A_{\mathrm{pair}}=(M-2)I+\mathbf 1\mathbf 1^\top$.
The unconstrained zero-covariance least-squares estimator satisfies
\begin{equation}
  \lVert\widehat{\mathbf r}_b-\mathbf r_b\rVert_2
  \leq \frac{2}{\sqrt{M-2}}\lVert\mathbf c_b\rVert_2.
  \label{eq:cov_bound}
\end{equation}
\end{theorem}
\begin{proof}
Equation~\eqref{eq:pair_system} gives
$\widehat{\mathbf r}_b-\mathbf r_b=-2A_{\mathrm{pair}}^+\mathbf c_b$.
The stated Gram matrix has eigenvalue $M-2$ orthogonal to $\mathbf 1$ and
$2(M-1)$ along $\mathbf 1$, so
$\lVert A_{\mathrm{pair}}^+\rVert_2=1/\sqrt{M-2}$.
\end{proof}

The theorem separates conditioning from systematic shared error. A larger,
more diverse committee improves the $1/\sqrt{M-2}$ factor, but adding many
correlated checkpoints does not make $\mathbf c_b$ vanish.

\subsection{Shift-conditioned covariance transport}

Target covariance cannot be measured without $Y_T$, so \method estimates it
from labeled source simulations. We construct pseudo-target graphs
\begin{equation}
  \widetilde{\mathcal G}_S^{(k)}
  =\mathcal T_{\xi_k}(\mathcal G_S,Y_S)
\end{equation}
that perturb features, topology, homophily, class prior, and conditional
structure. Frozen candidates run inference on each graph. Since $Y_S$ remains
known, we measure its band risks $r_{m,b}^{(k)}$ and error covariances
$c_{mn,b}^{(k)}$ without touching real target labels.

Each simulated graph and the real target is represented by an unlabeled shift
descriptor $\delta$ containing degree, spectral-moment, feature-smoothness,
embedding, and committee-prediction statistics. We match the target to the
simulation bank with simplex weights
\begin{equation}
  \alpha_T=\arg\min_{\alpha\in\Delta^K}
  \left\lVert\Sigma^{-1/2}
  \left(\sum_k\alpha_k\delta_k-\delta_T\right)\right\rVert_2^2
  +\lambda_\alpha\lVert\alpha\rVert_2^2,
  \label{eq:descriptor_match}
\end{equation}
where $\Delta^K=\{\alpha\geq0:\sum_k\alpha_k=1\}$ and descriptor variances
are floored. The transported prior risk and covariance are
\begin{equation}
  \widetilde r_{m,b}=\sum_k\alpha_{T,k}r_{m,b}^{(k)},
  \qquad
  \widehat c_{mn,b}=\sum_k\alpha_{T,k}c_{mn,b}^{(k)}.
  \label{eq:transport}
\end{equation}

\begin{corollary}[Transport error]
If Equation~\eqref{eq:pair_system} is corrected with
$\widehat{\mathbf c}_b$, its unconstrained least-squares solution satisfies
\begin{equation}
  \lVert\widehat{\mathbf r}_b-\mathbf r_b\rVert_2
  \leq\frac{2}{\sqrt{M-2}}
  \lVert\widehat{\mathbf c}_b-\mathbf c_b\rVert_2.
  \label{eq:transport_bound}
\end{equation}
\end{corollary}

Equation~\eqref{eq:transport_bound} states the empirical requirement plainly:
the unlabeled descriptor geometry must predict the covariance components that
matter for recovery. The assumption is testable and can fail under an unseen
shift family; Section~\ref{sec:experiments} reports this limitation rather than
hiding it behind aggregate selection performance.

\subsection{Practical recovery and optional uncertainty}

The practical solver combines transported covariance with reliability weights
and a support-aware prior:
\begin{align}
  \widehat{\mathbf r}_b
  =\arg\min_{\mathbf r_b\geq0}\;&
  \left\lVert W_b^{1/2}
  \left[A_{\mathrm{pair}}\mathbf r_b-
  (\mathbf d_b+2\widehat{\mathbf c}_b)\right]
  \right\rVert_2^2 \\
  &+\lambda_r\lVert\mathbf r_b\rVert_2^2
  +\lambda_p\lVert\mathbf r_b-\widetilde{\mathbf r}_b\rVert_2^2.
  \label{eq:practical_recovery}
\end{align}
Pairs with high transported error correlation receive less weight, and the
prior is disabled when the target descriptor falls outside the frozen
source-simulation support. The frozen implementation additionally winsorizes
extreme pair observations and performs one Tukey/endpoint-reliability refit.
We treat this as a solver detail: Table~\ref{tab:core_ablation} shows that it
does not explain the aggregate gain.

Optionally, bootstrap resampling of simulated shifts yields risk estimates
$\widehat R_T^{(q)}(h_m)$ and uncertainty
$u_m=\operatorname{Std}_q[\widehat R_T^{(q)}(h_m)]$. The final selector uses
\begin{equation}
  \widehat m=\arg\min_m
  \left[\frac{1}{2}\sum_b\widehat r_{m,b}+\beta u_m\right].
  \label{eq:ucb}
\end{equation}
All resampling operates on source-simulated calibration data, so the score
remains target-label-free. Its empirical contribution is reported separately
from covariance transport in Table~\ref{tab:core_ablation}.

\paragraph{Attribution boundary and complexity.}
With an exact tight frame, identical linear recovery in every band, and no
band-dependent calibration, summing spectral estimates equals applying the
same recovery operator to global disagreement. Therefore spectral value must
be attributed to band-conditioned covariance, weights, constraints, or
uncertainty---not to decomposition alone. Joint Chebyshev filtering costs
$O(BK|E_T|MC)$, while band-wise Gram matrices form all pair disagreements in
$O(BM^2n_TC)$. Recovery designs are cached by committee size; source-shift
calibration requires candidate inference but no GNN retraining.
